\documentclass[11pt]{article}
\usepackage[a4paper,margin=0.92in]{geometry}
\usepackage{amsmath,amssymb,amsthm}
\usepackage{mathrsfs}
\usepackage{enumitem}
\usepackage[hidelinks]{hyperref}

\newtheorem{theorem}{Theorem}
\newtheorem{lemma}{Lemma}
\newcommand{\Lz}{L^0(\mathcal F)}
\newcommand{\Lzp}{L^0_{+}(\mathcal F)}
\newcommand{\Lzpp}{L^0_{++}(\mathcal F)}
\newcommand{\Tc}{\mathcal T_c}
\newcommand{\Tel}{\mathcal T_{\varepsilon,\lambda}}
\newcommand{\Pcc}{\mathcal P_{cc}}
\newcommand{\ind}{\widetilde I}
\newcommand{\co}{\operatorname{Conv}_{L^0}}

\title{A random Schauder--Tychonoff fixed point theorem}
\author{FA/Banach run packet for arXiv:2104.11095}
\date{13 August 2026}

\begin{document}
\setlength{\emergencystretch}{3em}
\maketitle

\section*{The open problem}

Problem 5.1 of T. Guo, Y. Wang, H.-K. Xu, G. X. Yuan and G. Chen,
\emph{The noncompact Schauder fixed point theorem in random normed modules
and its applications}, arXiv:2104.11095, asks the following.  Let
$(E,\mathcal P)$ be a random locally convex module and let $G\subset E$ be
a stably compact $L^0$-convex set.  Must every $\sigma$-stable self-map of
$G$ which is either $\Tel(\mathcal P)$-continuous or
$\Tc(\mathcal P)$-continuous have a fixed point?

The source proves this for random normed modules.  A 2024 theorem of Q. Tu,
X. Mu and T. Guo proves the random Markov--Kakutani theorem in random
locally convex modules, but only for $L^0$-affine maps.  We give an
affirmative answer for the arbitrary nonlinear maps in Problem 5.1.

We use the standard notation of those papers.  Thus $\Pcc$ is the
$\sigma$-stable hull of $\mathcal P$: its members are countable
concatenations of finite maxima of seminorms in $\mathcal P$.  Stable
compactness may equivalently be taken with respect to
$\Tc(\Pcc)$.  As in the cited framework, modules are regular and all
equalities and inequalities are modulo null sets.

\section*{Full resolution}

\begin{theorem}[Random Schauder--Tychonoff theorem]\label{thm:main}
Let $(E,\mathcal P)$ be a random locally convex module over
$\mathbb K\in\{\mathbb R,\mathbb C\}$ with base $(\Omega,\mathcal F,P)$,
and let $G$ be a nonempty stably compact $L^0$-convex subset of $E$.
If $T:G\to G$ is $\sigma$-stable and is either
\begin{enumerate}[label=(\roman*),leftmargin=2.1em,nosep]
\item $\Tel(\mathcal P)$-continuous, or
\item $\Tc(\mathcal P)$-continuous,
\end{enumerate}
then $T$ has a fixed point.
\end{theorem}

The proof has two layers.  Stable compactness first produces a
piecewise-finite net for one seminorm.  The random Brouwer theorem then
gives an approximate fixed point for that seminorm.  Finally, stable
compactness is used a second time on the family of all closed
approximate-fixed-point sets.

\section*{Stable finite nets}

For $q\in\Pcc$, $r\in\Lzpp$ and $x\in G$, put
\[
 B_G(x;q,r)=\{y\in G:q(y-x)<r\text{ on }\Omega\}.
\]

\begin{lemma}[Piecewise-finite seminorm net]\label{lem:net}
For every $q\in\Pcc$ and $r\in\Lzpp$ there are a countable partition
$(A_n)$ of $\Omega$ and finite sets
\[
 D_n=\{x_{n1},\ldots,x_{nk_n}\}\subset G
\]
such that, for every $z\in G$ and every $n$,
\begin{equation}\label{eq:net}
 \ind_{A_n}\min_{1\le i\le k_n}q(z-x_{ni})
 <\ind_{A_n}r.
\end{equation}
\end{lemma}

\begin{proof}
The family
$\mathscr O=\{B_G(x;q,r):x\in G\}$ is a $\sigma$-stable family of
$\sigma$-stable $\Tc(\Pcc)$-open sets.  Indeed, locality of $q$ gives
\[
 \sum_m\ind_{C_m}B_G(x_m;q,r)
 =B_G\!\left(\sum_m\ind_{C_m}x_m;q,r\right)
\]
for every measurable partition $(C_m)$.  It covers $G$.  The stable
open-cover characterization of stable compactness therefore supplies a
stable finite subcover.  By definition, it has the form
\[
 \sum_n\ind_{A_n}\sigma(\mathscr O_n),
 \qquad
 \mathscr O_n=\{B_G(x_{ni};q,r):1\le i\le k_n\}.
\]
If $z\in G$, membership in this subcover gives, on $A_n$, a measurable
partition among the finitely many balls in $\mathscr O_n$.  Equivalently,
the pointwise minimum in \eqref{eq:net} is strictly smaller than $r$.
\end{proof}

\section*{One-seminorm approximate fixed points}

\begin{lemma}\label{lem:approx}
Under the hypotheses of Theorem \ref{thm:main}, for every $q\in\Pcc$ and
$r\in\Lzpp$ there exists $x\in G$ such that
\begin{equation}\label{eq:approx}
 q(x-Tx)<r\quad\text{on }\Omega.
\end{equation}
\end{lemma}

\begin{proof}
Choose $(A_n)$ and $D_n$ from Lemma \ref{lem:net}.  Fix $n$ and write
$k=k_n$, $x_i=x_{ni}$ and $A=A_n$.  On the standard random simplex
\[
 \Delta_k=\{\alpha=(\alpha_1,\ldots,\alpha_k)\in(\Lzp)^k:
                    \textstyle\sum_i\alpha_i=1\},
\]
define $h(\alpha)=\sum_i\alpha_i x_i\in G$.  For $z\in G$, set
\begin{align*}
 u_1(z)&=\ind_A\bigl(r-q(z-x_1)\bigr)_+ +\ind_{A^c},\\
 u_i(z)&=\ind_A\bigl(r-q(z-x_i)\bigr)_+,
       \qquad 2\le i\le k,
\end{align*}
and $U(z)=\sum_i u_i(z)$.  By \eqref{eq:net}, $U(z)>0$ on $A$;
by construction $U(z)=1$ on $A^c$.  Hence
\begin{equation}\label{eq:Phi}
 \Phi(\alpha)=
 \left(\frac{u_1(T(h(\alpha)))}{U(T(h(\alpha)))},\ldots,
       \frac{u_k(T(h(\alpha)))}{U(T(h(\alpha)))}\right)
\end{equation}
is a well-defined self-map of $\Delta_k$.

The map $\Phi$ is $\sigma$-stable.  It is also continuous in the topology
needed for the random Brouwer theorem.  In case (i),
$\Tel(\mathcal P)=\Tel(\Pcc)$, and $q$, $h$, positive part, addition,
multiplication and inversion on $L^0_{++}$ are continuous for the
corresponding $(\varepsilon,\lambda)$ topologies.  In case (ii), the 2024
continuity lemma says that a $\sigma$-stable
$\Tc(\mathcal P)$-continuous map is $\Tc(\Pcc)$-continuous; the same
operations are $\Tc$-continuous.  After restriction to the finite-rank
random simplex, either continuity hypothesis gives the random sequential
continuity used by the source's complete random Brouwer theorem.

Consequently, $\Phi$ has a fixed point $\alpha^n\in\Delta_k$.  Put
$y_n=h(\alpha^n)$.  On $A_n$, the fixed-point identity says
\[
 y_n=\sum_{i=1}^{k_n}
 \frac{(r-q(Ty_n-x_{ni}))_+}
      {\sum_j(r-q(Ty_n-x_{nj}))_+}\,x_{ni}.
\]
Only indices satisfying $q(Ty_n-x_{ni})<r$ receive positive weight.
The seminorm inequality therefore gives
\begin{equation}\label{eq:localapprox}
 \ind_{A_n}q(y_n-Ty_n)<\ind_{A_n}r.
\end{equation}
Finally concatenate $x=\sum_n\ind_{A_n}y_n$.  The set $G$ and the map
$T$ are $\sigma$-stable, while $q$ is local.  Thus $x\in G$ and
\eqref{eq:localapprox} concatenates to \eqref{eq:approx}.
\end{proof}

This is precisely the source paper's random Schauder projection, but run
on the finitely many centers available on each stable stratum and applied
to one seminorm at a time.  The coefficient-simplex formulation avoids
quotienting by $\ker q$; quotienting would be invalid because a nonlinear
$T$ need not descend to that quotient.

\section*{From approximate to exact}

For $q\in\Pcc$ and $r\in\Lzpp$, define
\[
 C(q,r)=\{x\in G:q(x-Tx)\le r\}.
\]
Lemma \ref{lem:approx} shows that every $C(q,r)$ is nonempty.  Each is
$\sigma$-stable.  It is also $\Tc(\Pcc)$-closed:
in case (ii) this follows from $\Tc(\Pcc)$-continuity; in case (i) it is
first $\Tel(\Pcc)$-closed, and the closure theorem for $\sigma$-stable sets
says that $\Tel(\mathcal P)$, $\Tc(\mathcal P)$ and $\Tc(\Pcc)$ have the
same $\sigma$-stable closed sets.

The family
\begin{equation}\label{eq:family}
 \mathscr C=\{C(q,r):q\in\Pcc,\ r\in\Lzpp\}
\end{equation}
is itself $\sigma$-stable.  To see this, let $(B_m)$ be a measurable
partition and let $C(q_m,r_m)\in\mathscr C$.  Then
\[
 q=\sum_m\ind_{B_m}q_m\in\Pcc,
 \qquad r=\sum_m\ind_{B_m}r_m\in\Lzpp,
\]
and locality of $T$ and the seminorms gives
\begin{equation}\label{eq:concatsets}
 \sum_m\ind_{B_m}C(q_m,r_m)=C(q,r).
\end{equation}
For the reverse inclusion in \eqref{eq:concatsets}, use a filler point in
each nonempty $C(q_m,r_m)$ off $B_m$ and concatenate; this turns the local
inequality for a point of $C(q,r)$ into a global member of each
$C(q_m,r_m)$.

The family has the finite-intersection property.  For
$C(q_i,r_i)$, $1\le i\le N$, take
\[
 q=\max_{1\le i\le N}q_i\in\Pcc,
 \qquad r=\min_{1\le i\le N}r_i\in\Lzpp.
\]
The assertion $q\in\Pcc$ follows by taking a common countable refinement
of the representations of the $q_i$; finite maxima remain finite maxima.
Then
\[
 \varnothing\ne C(q,r)\subseteq\bigcap_{i=1}^N C(q_i,r_i).
\]

Stable compactness and its closed-set finite-intersection characterization
now imply $\bigcap\mathscr C\ne\varnothing$.  Choose $x$ in this
intersection.  For every $p\in\mathcal P$ and every positive integer $m$,
\[
 p(x-Tx)\le \frac1m.
\]
Hence $p(x-Tx)=0$ for all $p\in\mathcal P$.  Since $\mathcal P$ separates
points, $x=Tx$.  This completes the proof of Theorem \ref{thm:main}.

\section*{Why the proof closes both continuity cases}

The weaker-topology case is not obtained by claiming that every
$\Tel$-continuous map on $G$ is globally $\Tc$-continuous.  Instead, the
only map to which Brouwer is applied is the finite-dimensional coefficient
map \eqref{eq:Phi}.  If $T$ is $\Tel$-continuous, then \eqref{eq:Phi} is
$\Tel$-continuous on a random Euclidean simplex, exactly a case covered by
the source's random Brouwer theorem.  If $T$ is $\Tc$-continuous, then
\eqref{eq:Phi} is $\Tc$-continuous there and the source's equivalence with
random sequential continuity applies.  The final closed-set argument uses
the equality of closures only for the $\sigma$-stable sets $C(q,r)$, where
that equality is valid.

\section*{Literature and novelty audit}

A bounded search through 13 August 2026 used the arXiv id, exact problem
language, title, authors, and combinations of ``stably compact'', ``random
locally convex module'', ``Schauder--Tychonoff'' and ``fixed point''.  The
closest later primary result found was the 2024 random Markov--Kakutani
theorem, which assumes $L^0$-affinity.  The 2025 random Kakutani theorem is
again restricted to random normed modules.  A March 2026 survey by several
of the same authors lists those results and describes the stable-compact
topological theory as still under development; it does not state the
nonlinear theorem above.  No posted full resolution of Problem 5.1 was
located.

\section*{References}

\begin{itemize}[leftmargin=2em]
\item T. Guo, Y. Wang, H.-K. Xu, G. X. Yuan and G. Chen,
\emph{The noncompact Schauder fixed point theorem in random normed modules
and its applications}, Math. Ann. 391 (2025), 3863--3911;
arXiv:2104.11095.  In particular: the complete random Brouwer theorem,
the random Schauder projection (Lemmas 3.7--3.9), and Problem 5.1.
\item Q. Tu, X. Mu and T. Guo, \emph{The random Markov--Kakutani fixed
point theorem in a random locally convex module}, New York J. Math. 30
(2024), 1196--1219.  In particular: Propositions 2.8, 2.14 and 2.15 and
Lemma 3.8.
\item Q. Tu, X. Mu, T. Guo, G. Yang and Y. Sun, \emph{The random Kakutani
fixed point theorem in random normed modules}, arXiv:2509.15649 (2025).
\item T. Guo, Q. Tu, X. Mu and Y. Sun, \emph{Survey of Metric fixed point
theory in random functional analysis}, arXiv:2603.27963 (2026).
\end{itemize}

\medskip
\noindent\textbf{Human-review recommendation.}
Review as a full resolution of Problem 5.1.  The proof-critical checks are:
(1) extracting \eqref{eq:net} from a stable finite subcover; (2) continuity
and $\sigma$-stability of \eqref{eq:Phi}; (3) the local fixed-point identity
and concatenation in Lemma \ref{lem:approx}; and (4) stable closedness and
the finite-intersection argument for \eqref{eq:family}.

\end{document}
